%% toc: P3 41 41.1
\section{Modules, abelian groups, and their definable structure}
\label{sec:p03-w03-p03-073-modules}

A module over a fixed ring is a structure whose first-order theory is
governed by one family of formulas, the positive primitive (pp) formulas.
These define subgroups, they are preserved by every homomorphism, and,
modulo the theory of modules, every formula reduces to a Boolean
combination of pp formulas and of sentences that count the indices of one
pp-definable subgroup in another. This section states that reduction
(pp-elimination), proves it in full, and draws its consequences:
elementary equivalence of modules is decided by indices of pp-definable
subgroups, every definable set of a module is a finite Boolean combination
of cosets of pp-definable subgroups, and every complete theory of modules
is stable in every cardinal \(\lambda\) with
\(\lambda^{\abs R+\aleph_0}=\lambda\). Abelian groups are the case
\(R=\mathbb Z\).

The inputs are a ring \(R\) (associative, with \(1\)), right
\(R\)-modules, and formulas of the language of \(R\)-modules. The
conclusions are the normal form of
\cref{thm:p03-w03-p03-073-pp-elimination}, the invariant criterion
\cref{thm:p03-w03-p03-073-elem-equiv}, the description of definable sets
\cref{thm:p03-w03-p03-073-definable-sets}, the type count
\cref{thm:p03-w03-p03-073-stability}, and the limits recorded in
\cref{thm:p03-w03-p03-073-limits}. The statements follow the survey of
Prest~\cite[\S\S2--7]{Prest2020ModelTheoryModules} and the summary in
Kuber~\cite[\S2.5]{Kuber2013GrothendieckModules}; the proofs below are
written out in the annex and use the compactness theorem
\cref{thm:fol-compactness}, the existence of saturated elementary
extensions \cref{thm:saturated-extension}, and the preservation of
positive existential formulas by homomorphisms
\cref{thm:p03-resplendence-positive-maps}. Pp formulas and pure
extensions for arbitrary signatures are introduced in
\cref{def:p03-resplendence-atomic} and
\cref{def:p03-resplendence-purity}; here they are specialized to modules.
Nilpotent groups of class two and further group examples follow in
\cref{sec:p03-w03-p03-074-nilpotent}.

\subsection{The language of modules and pp formulas}
\label{sec:p03-w03-p03-073-language}

\begin{hbconvention}{def:p03-w03-p03-073-language}{The language \(L_R\) and the theory \(T_R\)}
Fix a ring \(R\), associative with \(1\). The language \(L_R\) has a
binary function symbol \(+\), a constant \(0\), and for each
\(r\in R\) a unary function symbol for multiplication by \(r\); following
Prest~\cite[\S2]{Prest2020ModelTheoryModules} we write \(xr\) for its
application to \(x\). The theory \(T_R\) consists of the axioms of abelian
groups (with \(-x\) defined as \(x(-1)\)) and the right-module laws
\[
 (x+y)r=xr+yr,\qquad x(r+s)=xr+xs,\qquad x(rs)=(xr)s,\qquad x1=x
\]
for all \(r,s\in R\). Its models are the right \(R\)-modules, and an
\(L_R\)-homomorphism between models is the same thing as an \(R\)-linear
map. The cardinality of the set of \(L_R\)-formulas is
\(\kappa_R=\abs R+\aleph_0\). Unless said otherwise, ``module'' means
right \(R\)-module and formulas have no parameters. For a formula
\(\varphi(\bar x)\) with \(\bar x=(x_1,\ldots,x_n)\) and a module \(M\),
\(\varphi(M)=\{\bar a\in M^n:M\models\varphi(\bar a)\}\); for a formula
\(\varphi(\bar x,\bar y)\) and \(\bar b\) from \(M\),
\(\varphi(M,\bar b)=\{\bar a\in M^n:M\models\varphi(\bar a,\bar b)\}\).
Abelian groups are the \(\mathbb Z\)-modules. In \(L_{\mathbb Z}\)
multiplication by an integer is a function symbol; in the bare group
language \(\{+,0\}\) multiplication by a nonnegative integer is a term and
multiplication by a negative integer is definable through the additive
inverse, so translating each symbol by its definition carries
\(L_{\mathbb Z}\)-definable sets of an abelian group to
\(\{+,0\}\)-definable ones, and the converse holds trivially.
\index{module!language of}
\end{hbconvention}

\begin{hbdefinition}{def:p03-w03-p03-073-pp}{Pp formulas and pp-definable subgroups}
A \emph{positive primitive} (pp) formula of \(L_R\) is a formula
\(\exists\bar z\,(\theta_1\wedge\cdots\wedge\theta_k)\) with \(k\ge1\),
each \(\theta_j\) an atomic formula \(t=t'\), and \(\bar z\) a possibly
empty tuple of variables~\cite[\S3]{Prest2020ModelTheoryModules}. If
\(\varphi(\bar x)\) is pp, \(\varphi(M)\) is a \emph{pp-definable subgroup}
of \(M^n\) (\cref{thm:p03-w03-p03-073-pp-subgroups} shows it is a
subgroup). Two pp formulas \(\varphi(\bar x),\psi(\bar x)\) in the same
free variables form a \emph{pp pair}. Up to \(T_R\)-equivalence an atomic
formula is a homogeneous \(R\)-linear equation
\(\sum_i x_ir_i=0\), so a pp formula is an existentially quantified system
of homogeneous linear equations~\cite[\S3]{Prest2020ModelTheoryModules};
the proofs below do not need this normal form. Typical pp formulas in one
variable are the annihilation formula \(xr=0\) and the divisibility formula
\(r\mid x\), that is \(\exists y\,(x=yr)\).
\index{pp formula}\index{positive primitive formula}
\end{hbdefinition}

\begin{hblemma}{thm:p03-w03-p03-073-pp-subgroups}{proved=annex:p03-w03-p03-073-pp-subgroups}{Pp formulas define subgroups and cosets}
Let \(M\) be a module.
\begin{enumerate}
\item For every \(L_R\)-term \(t(x_1,\ldots,x_n)\) the map
      \(t^M:M^n\to M\) is a group homomorphism.
\item If \(\varphi(\bar x)\) is pp, then \(\varphi(M)\) is a subgroup of
      \(M^n\). In particular every pp sentence is true in every module.
\item The conjunction of two pp formulas is logically equivalent to a pp
      formula, \(\exists x_i\varphi\) is pp when \(\varphi\) is, and
      substituting terms (for instance \(0\)) for free variables of a pp
      formula yields a pp formula.
\item If \(\varphi(\bar x,\bar y)\) is pp and \(\bar b\) is a tuple from
      \(M\), then \(\varphi(M,\bar b)\) is either empty or a coset
      \(\bar a_0+\varphi(M,\bar 0)\) of the pp-definable subgroup
      \(\varphi(M,\bar 0)\).
\item If \(f:M\to N\) is \(R\)-linear and \(M\models\varphi(\bar a)\)
      with \(\varphi\) pp, then \(N\models\varphi(f\bar a)\)
      \cite[Lemma~3.1]{Prest2020ModelTheoryModules}.
\end{enumerate}
\end{hblemma}

\subsection{Invariants}
\label{sec:p03-w03-p03-073-invariants}

\begin{hbdefinition}{def:p03-w03-p03-073-invariants}{Invariants of a module and invariant sentences}
For a module \(M\) and a pp pair \(\varphi(\bar x),\psi(\bar x)\), the
\emph{invariant} \(\operatorname{Inv}(M;\varphi,\psi)\) is the index
\([\varphi(M):\varphi(M)\cap\psi(M)]\) if this is finite, and \(\infty\)
otherwise; this is Kuber's definition for pp-definable subgroups
\cite[Definition~2.5.4]{Kuber2013GrothendieckModules}. When \(\psi\)
implies \(\varphi\), Prest writes the same number as
\(\abs{\varphi(M)/\psi(M)}\)~\cite[\S3]{Prest2020ModelTheoryModules} and,
for the theory, \(\abs{\varphi/\psi}\)~\cite[\S7]{Prest2020ModelTheoryModules}. For an integer
\(k\ge1\) the \emph{invariant sentence} \(\operatorname{Inv}(\varphi,\psi)\ge k\)
is
\[
 \exists\bar x_1\cdots\exists\bar x_k\Bigl(\bigwedge_{i=1}^k\varphi(\bar x_i)
 \wedge\bigwedge_{1\le i<l\le k}\neg(\varphi\wedge\psi)(\bar x_i-\bar x_l)\Bigr),
\]
where \(\bar x_i-\bar x_l\) is the tuple of terms
\(x_{i,1}+x_{l,1}(-1),\ldots\). An \emph{invariants statement} is a
Boolean combination of invariant sentences
\cite[\S3]{Prest2020ModelTheoryModules}.
\index{invariant sentence}\index{module!invariants of}
\end{hbdefinition}

\begin{hblemma}{thm:p03-w03-p03-073-invariant-sentences}{proved=annex:p03-w03-p03-073-invariant-sentences}{What invariant sentences express}
For every module \(M\), pp pair \(\varphi,\psi\) and \(k\ge1\),
\(M\models\operatorname{Inv}(\varphi,\psi)\ge k\) iff
\(\operatorname{Inv}(M;\varphi,\psi)\ge k\). Consequently two modules
satisfy the same invariant sentences iff
\(\operatorname{Inv}(M;\varphi,\psi)=\operatorname{Inv}(N;\varphi,\psi)\)
for every pp pair, both values being read in
\(\{1,2,\ldots\}\cup\{\infty\}\).
\end{hblemma}

\subsection{Pp-elimination}
\label{sec:p03-w03-p03-073-pp-elimination}

The proof of pp-elimination rests on a covering lemma for cosets, which
Kuber states as Neumann's lemma and attributes to Prest's 1988 book
\cite[Lemma~2.5.8]{Kuber2013GrothendieckModules}.

\begin{hblemma}{thm:p03-w03-p03-073-neumann}{proved=annex:p03-w03-p03-073-neumann}{Neumann's lemma on coverings by cosets}
Let \(K\) be an abelian group, \(H,G_1,\ldots,G_m\) subgroups of \(K\)
with \(m\ge1\), and \(c,c_1,\ldots,c_m\in K\). If
\(c+H\subseteq\bigcup_{i=1}^m(c_i+G_i)\), then
\[
 c+H\subseteq\bigcup_{i\in I_0}(c_i+G_i),\qquad
 I_0=\{i:[H:H\cap G_i]<\infty\}.
\]
In particular, if a group is covered by finitely many cosets of its
subgroups, one of those subgroups has finite index.
\end{hblemma}

\begin{hblemma}{thm:p03-w03-p03-073-extension}{proved=annex:p03-w03-p03-073-extension}{One-variable extension step}
Let \(M,N\) be modules satisfying the same invariant sentences, and let
\(\bar a\in M^n\), \(\bar b\in N^n\) satisfy the same pp formulas
\(\varphi(\bar x)\). Let \(\varphi(\bar x,y)\) and
\(\psi_1(\bar x,y),\ldots,\psi_m(\bar x,y)\) be pp, \(m\ge0\). If some
\(e\in M\) satisfies
\(\varphi(\bar a,e)\wedge\neg\psi_1(\bar a,e)\wedge\cdots\wedge\neg\psi_m(\bar a,e)\),
then some \(f\in N\) satisfies
\(\varphi(\bar b,f)\wedge\neg\psi_1(\bar b,f)\wedge\cdots\wedge\neg\psi_m(\bar b,f)\).
\end{hblemma}

\begin{hblemma}{thm:p03-w03-p03-073-separation}{proved=annex:p03-w03-p03-073-separation}{Separation criterion for Boolean normal forms}
Assume the Axiom of Choice. Let \(T\) be an \(\Lang\)-theory, \(\Delta\) a
nonempty set of \(\Lang\)-formulas with free variables among
\(\bar x\), and \(\chi(\bar x)\) an \(\Lang\)-formula. Suppose that whenever
\(M,N\models T\), \(\bar a\in M^n\), \(\bar b\in N^n\), and \(\bar a\)
and \(\bar b\) satisfy the same members of \(\Delta\), then
\(M\models\chi(\bar a)\) iff \(N\models\chi(\bar b)\). Then there is a
Boolean combination \(\beta(\bar x)\) of members of \(\Delta\) with
\(T\models\forall\bar x\,(\chi\leftrightarrow\beta)\).
\end{hblemma}

\begin{hbtheorem}{thm:p03-w03-p03-073-pp-elimination}{proved=annex:p03-w03-p03-073-pp-elimination}{Pp-elimination of quantifiers}
Assume the Axiom of Choice. For every \(L_R\)-formula
\(\chi(\bar x)\) with free variables among \(\bar x=(x_1,\ldots,x_n)\),
\(n\ge0\), there is a Boolean combination \(\beta(\bar x)\) of pp formulas
with free variables among \(\bar x\) and of invariant sentences such that
\[
 T_R\models\forall\bar x\,\bigl(\chi(\bar x)\leftrightarrow\beta(\bar x)\bigr).
\]
Moreover, if \(M,N\) are modules with the same invariants and
\(\bar a\in M^n\), \(\bar b\in N^n\) satisfy the same pp formulas, then
\(\bar a\) and \(\bar b\) satisfy the same \(L_R\)-formulas.
\end{hbtheorem}

Prest states the theorem as the equivalence of every formula, modulo the
theory of right \(R\)-modules, with the conjunction of an invariants
statement and a Boolean combination of pp formulas
\cite[Theorem~4.1]{Prest2020ModelTheoryModules}; Kuber states it with a
finite disjunction of conjunctions of one pp formula and finitely many
negated pp formulas, together with an invariants statement, and attributes
it to Baur's 1976 paper
\cite[Theorem~2.5.5]{Kuber2013GrothendieckModules}. The Boolean combination
above contains both forms after rearrangement into disjunctive normal form.

\begin{hbcorollary}{thm:p03-w03-p03-073-elem-equiv}{proved=annex:p03-w03-p03-073-elem-equiv}{Invariants characterize elementary equivalence}
Assume the Axiom of Choice. Two modules \(M,N\) are elementarily
equivalent iff \(\operatorname{Inv}(M;\varphi,\psi)=\operatorname{Inv}(N;\varphi,\psi)\)
for every pp pair \(\varphi,\psi\). Hence every complete theory of
modules is axiomatized by \(T_R\) together with the invariant sentences
and negated invariant sentences it contains~\cite[\S7]{Prest2020ModelTheoryModules}.
\end{hbcorollary}

\begin{hbcorollary}{thm:p03-w03-p03-073-definable-sets}{proved=annex:p03-w03-p03-073-definable-sets}{Definable sets of a module}
Assume the Axiom of Choice. Let \(M\) be a module and \(n\ge1\). Every
subset of \(M^n\) definable with parameters from \(M\) is a finite Boolean
combination of sets \(\varphi(M,\bar b)\) with \(\varphi(\bar x,\bar y)\)
pp, each of which is empty or a coset of the pp-definable subgroup
\(\varphi(M,\bar 0)\) of \(M^n\)
\cite[after Theorem~4.1]{Prest2020ModelTheoryModules},
\cite[Theorem~2.5.6]{Kuber2013GrothendieckModules}.
\end{hbcorollary}

\subsection{Direct sums, products and stability}
\label{sec:p03-w03-p03-073-sums}

\begin{hblemma}{thm:p03-w03-p03-073-sums}{proved=annex:p03-w03-p03-073-sums}{Pp formulas commute with direct sums and products}
Let \((M_i)_{i\in I}\) be modules and \(\varphi(\bar x)\) pp with \(n\)
free variables. Assume the Axiom of Choice for the arbitrary-product
assertion; the direct-sum assertion and the two-module index formula do not
need this assumption. Identify \((\bigoplus_iM_i)^n\) with
\(\bigoplus_i M_i^n\) and \((\prod_iM_i)^n\) with \(\prod_iM_i^n\). Then
\[
 \varphi\Bigl(\bigoplus_iM_i\Bigr)=\bigoplus_i\varphi(M_i),\qquad
 \varphi\Bigl(\prod_iM_i\Bigr)=\prod_i\varphi(M_i).
\]
For two modules and a pp pair,
\(\operatorname{Inv}(M\oplus N;\varphi,\psi)
=\operatorname{Inv}(M;\varphi,\psi)\cdot\operatorname{Inv}(N;\varphi,\psi)\),
with \(\infty\cdot k=k\cdot\infty=\infty\) for \(k\ge1\)
\cite[\S7]{Prest2020ModelTheoryModules}.
\end{hblemma}

\begin{hbtheorem}{thm:p03-w03-p03-073-stability}{proved=annex:p03-w03-p03-073-stability}{Complete theories of modules are stable}
Assume the Axiom of Choice. Let \(M\) be a module, \(T=\operatorname{Th}(M)\)
and \(\kappa_R=\abs R+\aleph_0\). For every infinite cardinal \(\lambda\)
with \(\lambda^{\kappa_R}=\lambda\), \(T\) is \(\lambda\)-stable in the
sense of \cref{def:p03-stability-spectrum}: for every \(N\models T\) and
\(A\subseteq\dom N\) with \(\abs A\le\lambda\), at most \(\lambda\)
complete \(1\)-types over \(A\) are consistent with the elementary diagram
of \(N\) restricted to \(A\). In particular \(T\) is
\(2^{\kappa_R}\)-stable.
\end{hbtheorem}

Prest records that every module is stable and derives it from the same
count of pp-types, bounding the number of pp-types over \(A\) by
\(\abs A^{\kappa_R}\)~\cite[\S6 and Theorem~6.1(i)]{Prest2020ModelTheoryModules}.
That theorem also characterizes superstability and total transcendence by
chain conditions on pp-definable subgroups; those parts are not proved here.

\subsection{Procedure, examples and limits}
\label{sec:p03-w03-p03-073-procedure}

\begin{hbmethod}{def:p03-w03-p03-073-method}{Comparing two modules by their invariants}
The invariant criterion compares modules \(M\) and \(N\) as follows:
\begin{enumerate}
\item determine, for each pp pair \(\varphi,\psi\), the subgroups
      \(\varphi(M)\), \((\varphi\wedge\psi)(M)\) and their counterparts
      in \(N\);
\item compare \(\operatorname{Inv}(M;\varphi,\psi)\) with
      \(\operatorname{Inv}(N;\varphi,\psi)\);
\item conclude \(M\equiv N\) iff all comparisons agree
      (\cref{thm:p03-w03-p03-073-elem-equiv}).
\end{enumerate}
Side conditions: both structures are modules over the same ring \(R\) in
the same language \(L_R\); the comparison ranges over all pp pairs in all
finite numbers of free variables, not only one-variable pairs, unless a
separate result reduces the pairs needed; and only finiteness or the exact
finite value of each index matters, never the size of an infinite index.
For direct sums, \cref{thm:p03-w03-p03-073-sums} reduces step~1 to the
summands. This is a structural criterion, not a terminating algorithm:
there are generally infinitely many pp pairs, and their indices need not
be effectively available. A decision procedure requires an additional
effective reduction of these comparisons.
\end{hbmethod}

\begin{hbcorollary}{thm:p03-w03-p03-073-vector-spaces}{proved=annex:p03-w03-p03-073-vector-spaces}{Nonzero rational vector spaces are elementarily equivalent abelian groups}
Assume the Axiom of Choice. Every nonzero \(\mathbb Q\)-vector space, as an
abelian group in \(L_{\mathbb Z}\), is elementarily equivalent to
\(\mathbb Q\). In particular \(\mathbb Q\equiv\mathbb Q\oplus\mathbb Q\).
\end{hbcorollary}

\begin{hbexample}{def:p03-w03-p03-073-example-finite}{Reading invariants of two groups of order four}
Take \(R=\mathbb Z\), \(M=\mathbb Z/4\) and
\(N=\mathbb Z/2\oplus\mathbb Z/2\). For the pp pair \(\varphi(x)\) given
by \(x=x\) and \(\psi(x)\) given by \(x2=0\), one has
\(\psi(M)=\{0,2\}\) and \(\psi(N)=N\), so
\(\operatorname{Inv}(M;\varphi,\psi)=2\) and
\(\operatorname{Inv}(N;\varphi,\psi)=1\). By
\cref{thm:p03-w03-p03-073-invariant-sentences} the invariant sentence
\(\operatorname{Inv}(\varphi,\psi)\ge2\) holds in \(M\) and fails in
\(N\), so \(M\not\equiv N\), although both groups have four elements. The
same pair computed on \(N\) through \cref{thm:p03-w03-p03-073-sums} gives
\(1\cdot1=1\), since \(\psi(\mathbb Z/2)=\mathbb Z/2\). For the
divisibility formula \(\delta(x)\), namely \(\exists y\,(x=y2)\),
\(\delta(M)=\{0,2\}=\psi(M)\) while \(\delta(N)=\{0\}\), so the pp pair
\(\delta,\psi\) separates in the other direction:
\(\operatorname{Inv}(M;\psi,\delta)=1\) and
\(\operatorname{Inv}(N;\psi,\delta)=4\).
\end{hbexample}

\begin{hbtheorem}{thm:p03-w03-p03-073-limits}{proved=annex:p03-w03-p03-073-limits}{Invariant sentences are needed, and pp-elimination is not quantifier elimination}
\begin{enumerate}
\item The sentence \(\exists y\,\neg(y=0)\) is not equivalent modulo
      \(T_R\), for any nonzero ring \(R\), to a Boolean combination of pp
      sentences.
\item In the abelian group \(\mathbb Z\) the pp-definable set
      \(2\mathbb Z=\delta(\mathbb Z)\), with \(\delta(x)\) the formula
      \(\exists y\,(x=y2)\), is not defined by any quantifier-free
      \(L_{\mathbb Z}\)-formula without parameters. Hence
      \(\operatorname{Th}(\mathbb Z)\) does not admit elimination of
      quantifiers in \(L_{\mathbb Z}\).
\end{enumerate}
\end{hbtheorem}

\begin{hbwarning}{def:p03-w03-p03-073-warning}{What pp-elimination does not give}
Pp-elimination is a relative normal form: it removes quantifiers only down
to pp formulas, which themselves contain existential quantifiers, and it
keeps invariant sentences, which cannot be dropped
(\cref{thm:p03-w03-p03-073-limits}). Complete elimination of quantifiers
for every complete theory of \(R\)-modules holds exactly when \(R\) is von
Neumann regular, according to Prest~\cite[Theorem~5.1]{Prest2020ModelTheoryModules};
that equivalence is not proved here, and \(\mathbb Z\), which is not von
Neumann regular, illustrates the failure in
\cref{thm:p03-w03-p03-073-limits}. The proofs here use compactness and
saturated elementary extensions and therefore the Axiom of Choice; no
claim is made about a choice-free derivation. The criterion
\cref{thm:p03-w03-p03-073-elem-equiv} is not by itself a decision
procedure. For abelian groups, \cref{thm:p03-ab-unary-normal} supplies
the effective integer form, \cref{thm:p03-ab-index-calculation} computes
the finite invariant conditions, and \cref{thm:p03-ab-realization}
decides their realizability. Together they prove Szmielew's theorem in
\cref{thm:p03-abelian-decidable}. This decides the common theory of all
abelian groups, not the complete theory of every individual group;
\cref{thm:p03-ab-individual-boundary} gives a counterexample to that
stronger conclusion.
\end{hbwarning}

\subsection{Relationships and sources}
\label{sec:p03-w03-p03-073-sources}

The pp formulas of this section are the positive primitive formulas of
\cref{def:p03-resplendence-atomic} for the signature \(L_R\);
an embedding of modules is pure in the sense of
\cref{def:p03-resplendence-purity} exactly when every pp formula with
parameters from the smaller module that holds in the larger one already
holds in the smaller one, and Prest takes this notion as the replacement
for elementary embeddings in the model theory of modules
\cite[\S2]{Prest2020ModelTheoryModules}. The extension step
\cref{thm:p03-w03-p03-073-extension} is the module form of the
back-and-forth arguments used for saturated models in
\cref{thm:saturation-homogeneity}. The stability count
\cref{thm:p03-w03-p03-073-stability} supplies examples of stable theories
for \cref{def:p03-stability-spectrum} that are not \(\omega\)-stable in
general; Prest notes that \(\mathbb Z\) as a module over itself is
superstable but not totally transcendental
\cite[\S6]{Prest2020ModelTheoryModules}.

Sources read for this section: Prest's survey \emph{Model Theory of
Modules}, version dated 4 September 2020 from Prest's web page,
Sections~1--7 (definitions, Lemma~3.1, Theorem~4.1, Theorem~5.1,
Theorem~6.1)~\cite{Prest2020ModelTheoryModules}; Kuber's
\emph{Grothendieck Rings of Theories of Modules}, arXiv version~1 of
18 February 2013, Section~2.5~\cite{Kuber2013GrothendieckModules}. The proofs in the annex are written
for this handbook and are not transcriptions of either source.
